\documentclass[11pt]{article}
\usepackage[margin=1in]{geometry}
\usepackage{amsmath,amssymb,amsthm,booktabs,enumitem,hyperref,natbib}
\hypersetup{
 colorlinks=true,
 linkcolor=blue,
 citecolor=blue,
 urlcolor=blue,
 pdftitle={Boundary partitions at the leximin core},
 pdfauthor={Maxwell Flanagan}
}

\newtheorem{theorem}{Theorem}
\newtheorem{lemma}[theorem]{Lemma}
\newtheorem{proposition}[theorem]{Proposition}
\newtheorem{corollary}[theorem]{Corollary}
\newtheorem{remark}[theorem]{Remark}
\newtheorem{definition}[theorem]{Definition}
\newtheorem{example}[theorem]{Example}
\newcommand{\core}{\operatorname{Core}}
\newcommand{\Lmin}{\operatorname{Lmin}}
\newcommand{\one}{\mathbf 1}
\newcommand{\PART}{\mathrm{BP}}
\newcommand{\TODO}[1]{\textbf{[TODO: #1]}}

\title{Boundary partitions at the leximin core:\\
a two-parameter characterization}
\author{Maxwell Flanagan}
\date{Draft of August 8, 2026}

\begin{document}
\maketitle

\begin{abstract}
At a slack strict payoff boundary of the leximin core allocation of a
singleton-normalized exact game, the two richer players may each carry a tight
coalition, and those two coalitions may partition the player set.  We
characterize exactly when: writing $p$ and $q$ for the number of poorer players
on the two sides, such a \emph{boundary partition} exists if and only if
$p,q\ge1$ and $p+q\ge4$.  Both directions are proved in closed form.  The
obstruction and the construction turn out to be the same inequality read in
opposite directions: a tight coalition that contains the boundary and omits one
side's poorer players costs two spare players on the other side.  The
characterization settles a residual question about zero-payoff configurations
on seven players, in the affirmative rather than, as we had conjectured, the
negative.  Two methodological ingredients are of independent use: exactness of
a game is an existential linear condition, needing no semi-balanced catalogue
and no vertex enumeration; and, at a fixed payoff stratification, leximin
optimality is a Farkas certificate supported on the tight coalitions, hence
quantifier-free.
\end{abstract}

\section{Introduction}

Let $v$ be a transferable-utility game on a finite player set $N$ that is
\emph{exact}, meaning $v(S)=\min_{x\in\core(v)}x(S)$ for every coalition, and
\emph{singleton-normalized}, meaning $v(\{i\})=0$ for every player.  Let
$\phi=\Lmin(v)$ be its leximin core allocation, and let $N=L\mathbin{\dot\cup}H$
be a strict payoff boundary: $H$ is a union of the richest payoff strata of
$\phi$ and $\max_{i\in L}\phi_i<\min_{j\in H}\phi_j$.  The boundary is
\emph{slack} when $v(H)<\phi(H)$.  A companion paper \citep{Flanagan2026phase}
determines exactly which pairs $(|L|,|H|)$ admit a slack boundary.

This note answers a finer question at the smallest interesting richer side,
$|H|=2$.  Write $H=\{f,g\}$.

\begin{definition}[boundary partition]\label{def:bp}
A \emph{boundary partition of type $(p,q)$} at a slack strict boundary
$L\mathbin{\dot\cup}H$ with $H=\{f,g\}$ is a pair of $\phi$-tight coalitions
\[
 S_1=P\cup\{f\},\qquad S_2=Q\cup\{g\},\qquad
 P\mathbin{\dot\cup}Q=L,\quad p=|P|,\ q=|Q|,
\]
so that $S_1$ and $S_2$ are disjoint, cover $N$, and meet $H$ in one player
each.  We write $\PART(p,q)$ for the type, and $m=p+q=|L|$, $n=m+2=|N|$.
\end{definition}

The configuration is the extreme case of the failure of exact partition: not
only is the richer suffix $H$ not tight, but the two richer players separately
anchor a partition of $N$ into tight blocks.  Our main result is that this is
possible exactly on a two-parameter staircase.

\begin{theorem}[characterization]\label{thm:main}
A boundary partition of type $(p,q)$ exists if and only if
\[
 p\ge1,\qquad q\ge1,\qquad p+q\ge4 .
\]
\end{theorem}

Necessity is Theorem~\ref{thm:necessity}; sufficiency is
Theorem~\ref{thm:sufficiency}, by two explicit families of games.  Neither
direction uses a solver.

Three features of the proof are worth advertising.

\paragraph{One inequality, read twice.}
Everything turns on a single consequence of the uncrossing budget for tight
coalitions (Lemma~\ref{lem:cross}): if $T$ is $\phi$-tight, contains $H$, and
omits every player of $P$, then
\begin{equation}\label{eq:key}
 |Q\setminus T|\ \ge\ 2 .
\end{equation}
Read as a prohibition, \eqref{eq:key} kills every type with $p+q\le3$.  Read as
a design constraint, it dictates the two realizing families and their ranges.
The threshold is therefore not two facts that happen to meet at $m=4$; it is
one inequality that first becomes satisfiable there.

\paragraph{A three-line normal form.}
Because $S_1$ and $S_2$ are complementary tight coalitions, \emph{every} core
point is tight on both (Lemma~\ref{lem:rigid}).  The core therefore collapses
to a polytope in $|L|$ variables, and the $2^{n}$ coalition functionals
collapse onto a short explicit list whose facet normals can never separate two
players lying in the same block (Lemma~\ref{lem:normal}).  That last sentence
is the combinatorial content of the configuration and is what the necessity
proof consumes.

\paragraph{A correction.}
An earlier round of this program conjectured that a boundary partition requires
$\max(p,q)\ge3$, on the strength of a census in which type $\PART(2,2)$ never
appeared among some twenty thousand sampled boundaries.  That conjecture is
false: $\PART(2,2)$ is realizable, and Section~\ref{sec:zp} exhibits a
six-player game carrying it.  The correct threshold is on $m=p+q$, not on
$\max(p,q)$.  We record this because the failure mode is a general hazard:
absence from a random census is evidence about sampling, not about
realizability, and the same caution is raised in \citep{Flanagan2026phase}
for a different configuration.  The consequence for the seven-player
zero-payoff question of Section~\ref{sec:zp} is that it resolves in the
opposite direction from the one we expected.

\paragraph{Prior art.}\label{sec:intro-priorart}
\TODO{prior-art positioning; the lead runs the citation verification pass and
no reference below has been invented.}  Citations are needed for: the
egalitarian core and its transfer-blocking characterization, which is where the
protector Lemma~\ref{lem:prot} sits in the literature; exact games and large
cores, for the standing hypotheses of Section~\ref{sec:prelim}; Dutta--Ray
egalitarianism and Lorenz domination, which is the context in which exact
partition at $\Lmin$ is normally discussed; and any prior appearance of an
existential characterization of exactness in the spirit of
Theorem~\ref{thm:exact}, which we believe to be folklore but have not traced.
The companion paper \citep{Flanagan2026phase} carries the corresponding
positioning for the phase diagram and is the right place to start.

\medskip
Sections~\ref{sec:prelim} and~\ref{sec:descent} fix notation and record the
descent that reduces zero-payoff configurations to boundary partitions.
Section~\ref{sec:methods} isolates the two encoding theorems.
Section~\ref{sec:frame} sets up the normal form, Section~\ref{sec:nec} proves
necessity, Section~\ref{sec:suf} gives the families.  Section~\ref{sec:frontier}
describes what happens exactly at the threshold, and Section~\ref{sec:zp} draws
the seven-player consequence.  Sections~\ref{sec:status}
and~\ref{sec:scope} state the status of every claim and the limits of the
result.

\section{Preliminaries}\label{sec:prelim}

All games are on a finite set $N$, exact, and singleton-normalized; core
allocations are then nonnegative.  For $\phi\in\core(v)$ write
$s_\phi(S)=\phi(S)-v(S)$ for the slack, and call $S$ \emph{$\phi$-tight} when
$s_\phi(S)=0$.  We use $\Lmin(v)$ for the leximin-maximal core allocation,
computed in the usual stages: at each stage one maximizes the common lower
bound on the payoffs of the not-yet-pinned players subject to the pinned ones
being fixed, and pins those players whose payoff cannot exceed that bound.  The
payoff strata of $\phi$ are its level sets, ordered $B_1\prec\cdots\prec B_k$
from poorest to richest, and a \emph{richer suffix} is a union
$B_{j+1}\cup\cdots\cup B_k$.

We use two facts from the companion paper.  The first is the uncrossing budget.

\begin{lemma}[crossing budget]\label{lem:cross}
Let $v$ be exact and singleton-normalized and $\phi\in\core(v)$.  If $S_1$ and
$S_2$ are $\phi$-tight then
\[
 s_\phi(S_1\cap S_2)\ \le\ \phi\bigl(N\setminus(S_1\cup S_2)\bigr).
\]
\end{lemma}

The second is the existence of protectors, which we restate with a direct proof
because it is used repeatedly.

\begin{lemma}[protectors]\label{lem:prot}
Let $\phi=\Lmin(v)$.  If $\phi_i<\phi_j$ then some proper $\phi$-tight
coalition contains $j$ and omits $i$.
\end{lemma}

\begin{proof}
Otherwise $d=e_i-e_j$ satisfies $d(T)=[i\in T]-[j\in T]\ge0$ for every
$\phi$-tight $T$, the value $-1$ occurring only when $j\in T$ and $i\notin T$.
Also $d(N)=0$, so $\phi+td\in\core(v)$ for all small $t>0$.  Choose
$0<t<(\phi_j-\phi_i)/2$.  The sorted payoff vector then changes only by
replacing $\phi_i$ with $\phi_i+t$ and $\phi_j$ with $\phi_j-t$, and
$\phi_j-t>\phi_i+t$, so no coordinate falls to or below $\phi_i$: the first
position at which the two sorted vectors differ is strictly larger for
$\phi+td$, contradicting leximin-maximality.
\end{proof}

Finally we use the standard lower-envelope construction.

\begin{lemma}[envelope]\label{lem:env}
Let $G$ be a finite family of nonnegative allocations with a common total $T$,
and set $v(S)=\min_{g\in G}g(S)$.  Then $v$ is exact and every $g\in G$ lies in
$\core(v)$; and $v$ is singleton-normalized as soon as every coordinate
vanishes in some row of $G$.
\end{lemma}

\section{Zero-payoff configurations and the descent}\label{sec:descent}

Boundary partitions arose from a neighbouring question.  Keep $H=\{f,g\}$ and
let $S_1,S_2$ be $\phi$-tight with $S_1\cap H=\{f\}$ and $S_2\cap H=\{g\}$.  A
\emph{zero-payoff configuration} of type $\mathrm{ZP}(p,q)$ has
$S_1\cap S_2=\{e\}$ for a single $e\in L$ and $S_1\cup S_2=N$, with
$p=|S_1\cap L\setminus\{e\}|$ and $q=|S_2\cap L\setminus\{e\}|$.  Boundary
partitions are the degenerate case $S_1\cap S_2=\emptyset$.

The two are linked by deletion and suspension of a null player.  We record the
chain; the arguments are short and are given in full in the archived proof file
cited in Section~\ref{sec:status}.

\begin{proposition}[descent]\label{prop:descent}
Let $v$ be exact and singleton-normalized with $\phi=\Lmin(v)$ and a slack
strict boundary $H=\{f,g\}$.
\begin{enumerate}[nosep,label=(\roman*)]
\item In a $\mathrm{ZP}(p,q)$ configuration, $\phi_e=0$.
\item If $\min_i\phi_i=0$ then a single player $j$ has $x_j=0$ for every
      $x\in\core(v)$, and $j$ is null.
\item Deleting a null player preserves exactness, normalization, the core, the
      leximin point, the tightness of every tight coalition, and the slack of
      every boundary avoiding it.
\item Conversely, adjoining a null player to a game with a boundary partition
      of type $(p,q)$ yields a $\mathrm{ZP}(p,q)$ configuration with one more
      poorer player.
\item Hence every $\mathrm{ZP}(p,q)$ configuration descends, one poorer player
      at a time, either to a boundary partition of type $(p,q)$ or to a
      zero-payoff configuration of type $(p-1,q)$ or $(p,q-1)$.
\end{enumerate}
\end{proposition}

Two immediate consequences will be used.  First, no boundary partition has
$q=0$: that would make the rich singleton $\{g\}$ tight, forcing
$\phi_g=v(\{g\})=0$ against strictness.  So $p,q\ge1$, which is the easy half
of Theorem~\ref{thm:main}.  Second, by (ii)--(iii) we may assume throughout
that no poorer player is null:

\begin{lemma}[no null poorer player]\label{lem:nonull}
If a boundary partition of type $(p,q)$ exists with $\min_{i\in L}\phi_i=0$,
then one of type $(p-1,q)$ or $(p,q-1)$ exists.  Consequently, in a boundary
partition with $m$ minimal, $\phi_i>0$ for every $i\in L$ and no singleton is
$\phi$-tight.
\end{lemma}

\section{Two encodings that decide}\label{sec:methods}

The results below are proved by hand, but the route to them ran through finite
checks, and the two encoding theorems that made those checks decisive are of
independent interest.  Both replace an apparently intractable condition by an
exactly equivalent linear one.

\begin{theorem}[exactness is existential]\label{thm:exact}
A game $v$ is exact if and only if for every coalition $S$ there is
$z^S\in\core(v)$ with $z^S(S)=v(S)$.
\end{theorem}

\begin{proof}
Immediate from $v(S)=\min_{x\in\core(v)}x(S)$: the minimum is attained exactly
when such a $z^S$ exists, and every core point satisfies $x(S)\ge v(S)$.
\end{proof}

The content is what this buys.  Exactness is usually imposed either by
enumerating core vertices, which is exponential, or through the facet
description of the cone of exact games by indecomposable min-semi-balanced
systems \citep{StudenyKratochvil2022,LohmannBormHerings2012}, which is
available only up to six players and whose sufficiency direction is delicate to
validate.  Theorem~\ref{thm:exact} instead adjoins one allocation vector per
coalition subject to linear constraints: no catalogue, no enumeration, no
completeness hypothesis, and no restriction on $n$.

The second theorem removes the quantifier in ``$\phi$ is the leximin point''.

\begin{theorem}[leximin is a Farkas certificate]\label{thm:farkas}
Let $v$ be a game, $\phi\in\core(v)$, and let $B_1\prec\cdots\prec B_k$ be the
payoff strata of $\phi$, with common values $\mu_1<\cdots<\mu_k$.  Write
$\mathcal T$ for the family of proper $\phi$-tight coalitions.  Then
$\phi=\Lmin(v)$ if and only if for every stage $t$ and every $i\in B_t$ there
are reals $y_A\ge0$ ($A$ proper), $\eta$, $\rho_l$ for
$l\in B_1\cup\cdots\cup B_{t-1}$, and $\sigma_c\ge0$ for $c\in B_t$, such that
\begin{equation}\label{eq:support}
 y_A=0\quad\text{for every }A\notin\mathcal T,
\end{equation}
and, for every player $j$,
\begin{equation}\label{eq:identity}
 -\sum_{A\ni j}y_A+\eta+[\,j\in B_{<t}\,]\rho_j-[\,j\in B_t\,]\sigma_j
  =[\,j=i\,].
\end{equation}
\end{theorem}

\begin{proof}
By the staged definition, $\phi=\Lmin(v)$ with these strata if and only if for
every $t$ and every $i\in B_t$,
\begin{equation}\label{eq:stage}
 \max\{x_i:\ x\in\core(v),\ x_l=\phi_l\ (l\in B_{<t}),\
      x_c\ge\mu_t\ (c\in B_{\ge t})\}=\phi_i .
\end{equation}
Necessity is the definition.  For sufficiency, \eqref{eq:stage} at stage $t$
forces the stage level to be exactly $\mu_t$ --- a core point with every
unpinned coordinate above $\mu_t$ would exceed $\mu_t$ at $i\in B_t$ --- and
forces the pinned set to be exactly $B_t$, since for $i$ in a later stratum
$\phi$ itself gives $x_i=\phi_i>\mu_t$; induct on $t$.

Now dualize \eqref{eq:stage}.  Its feasible set is a polyhedron containing
$\phi$, whose constraints are efficiency; $x(A)\ge v(A)$ for proper $A$, active
exactly on $\mathcal T$; the pins $x_l=\phi_l$; and $x_c\ge\mu_t$, active
exactly for $c\in B_t$, because $\phi_c>\mu_t$ for $c$ in a later stratum.
Nonnegativity is the singleton case of $x(A)\ge v(A)$.  The tangent cone at
$\phi$ is therefore
\[
 \mathcal D_t=\{d:\ d(N)=0,\ d(T)\ge0\ (T\in\mathcal T),\
   d_l=0\ (l\in B_{<t}),\ d_c\ge0\ (c\in B_t)\},
\]
and \eqref{eq:stage} says $d_i\le0$ on $\mathcal D_t$, i.e.\ $e_i$ lies in the
polar cone.  That polar is generated by $-\one_T$ for $T\in\mathcal T$, by
$-e_c$ for $c\in B_t$, and by the lines through $\one_N$ and through $e_l$ for
$l\in B_{<t}$; membership of $e_i$ in it is exactly \eqref{eq:identity} with
\eqref{eq:support}.  Conversely, given such multipliers, for any
$d\in\mathcal D_t$,
$d_i=-\sum_Ay_Ad(A)+\eta\,d(N)+\sum_l\rho_ld_l-\sum_c\sigma_cd_c\le0$,
each term being non-positive.
\end{proof}

\begin{remark}
The naive dual of \eqref{eq:stage} carries products $y_Av(A)$, which are
bilinear in the unknowns when the game is itself unknown.  Complementary
slackness at the optimum $\phi$ removes them: $y_A>0$ forces $A$ tight, so
$\sum_Ay_Av(A)=\sum_Ay_A\phi(A)$, and pairing \eqref{eq:identity} with $\phi$
reproduces the numerical half of the dual automatically.  What survives is
\eqref{eq:support}, a disjunction rather than a product.  This is why the
condition is decidable in linear arithmetic with Boolean structure.
\end{remark}

\section{The block-polytope normal form}\label{sec:frame}

Fix a boundary partition of type $(p,q)$ with the notation of
Definition~\ref{def:bp}.

\begin{lemma}[block rigidity]\label{lem:rigid}
Every $x\in\core(v)$ satisfies $x(S_1)=v(S_1)$ and $x(S_2)=v(S_2)$.  In
particular $v(S_1)+v(S_2)=v(N)$, and the core lies in a face of codimension
one.
\end{lemma}

\begin{proof}
Tightness at $\phi$ and efficiency give
$v(S_1)+v(S_2)=\phi(S_1)+\phi(S_2)=\phi(N)=v(N)$.  For $x\in\core(v)$,
$x(S_1)\ge v(S_1)$ and $x(S_2)\ge v(S_2)$ while
$x(S_1)+x(S_2)=x(N)=v(S_1)+v(S_2)$.  Two inequalities summing to an equality
are equalities.
\end{proof}

Lemma~\ref{lem:rigid} is the precise sense in which a boundary partition is
close to a direct sum: the core is confined to the product of the two block
faces.  It is not a direct sum --- the core is in general strictly inside that
product --- and the gap is exactly what makes the threshold $4$ rather than
something read off one block at a time.

Write $u_i=x_i-\phi_i$ for $i\in L$ and $U_P=\sum_{i\in P}u_i$,
$U_Q=\sum_{i\in Q}u_i$.

\begin{lemma}[normal form]\label{lem:normal}
Under Lemma~\ref{lem:rigid} the core is a polytope $K\subseteq\mathbb R^L$
containing the origin, with $x_f-\phi_f=-U_P$ and $x_g-\phi_g=-U_Q$.  For a
coalition $T$ put $\varepsilon_f=[f\in T]$ and $\varepsilon_g=[g\in T]$.  Then
the displacement functional $d(T)$ is the linear form on $\mathbb R^L$ with
coefficients
\[
 \nu_i=[i\in T]-\varepsilon_f\ (i\in P),\qquad
 \nu_i=[i\in T]-\varepsilon_g\ (i\in Q).
\]
Consequently, within each block the coefficients lie all in $\{0,1\}$ or all in
$\{-1,0\}$: \emph{no facet normal separates two players of the same block with
opposite signs}.  Among proper nonempty coalitions the zero form is attained
exactly at $T=S_1$ and $T=S_2$.
\end{lemma}

\begin{proof}
Solving $x(S_1)=v(S_1)$ for $x_f$ and $x(S_2)=v(S_2)$ for $x_g$ gives the
substitutions.  Then
$d(T)=\sum_{i\in T\cap L}u_i-\varepsilon_fU_P-\varepsilon_gU_Q$, and
collecting the coefficient of $u_i$ gives the display.  For $i\in P$ that
coefficient is $[i\in T]-\varepsilon_f$, which ranges over $\{0,1\}$ when
$\varepsilon_f=0$ and over $\{-1,0\}$ when $\varepsilon_f=1$; likewise on $Q$.
\end{proof}

The tangent cone of the core at $\phi$ is
$\mathcal D=\{d\in\mathbb R^L:\nu_T\cdot d\ge0$ for every $\phi$-tight $T\}$.

\begin{lemma}[improving directions]\label{lem:improve}
$\mathcal D\cap\mathbb R^L_{\ge0}=\{0\}$.  Consequently, for every nonempty
$S\subseteq L$ there is a $\phi$-tight $T$ with
\begin{equation}\label{eq:istar}
 |T\cap S|\ <\ \varepsilon_f\,|P\cap S|+\varepsilon_g\,|Q\cap S| .
\end{equation}
\end{lemma}

\begin{proof}
Let $d\in\mathcal D$ be nonnegative on $L$ and nonzero.  The boundary is
strict, so every stratum of $L$ precedes every stratum meeting $H$; let $B_j$
be the first stratum on which $d$ is nonzero.  Since $d\ne0$ somewhere on $L$,
$B_j$ is a stratum of $L$, so $d|_{B_j}\ge0$ and $d|_{B_j}\ne0$.  Then
$\phi+td$ lies in the core for small $t>0$ and raises the sorted payoff vector
at the first position where it changes, contradicting $\phi=\Lmin(v)$.  For the
second assertion, apply the first to $d=\one_S$, which is nonnegative and
nonzero, so $\one_S\notin\mathcal D$; by Lemma~\ref{lem:normal},
$\nu_T\cdot\one_S=|T\cap S|-\varepsilon_f|P\cap S|-\varepsilon_g|Q\cap S|$ for
every $T$, and some $\phi$-tight $T$ makes it negative.
\end{proof}

\begin{remark}\label{rem:noconverse}
We claim only the stated direction.  Condition \eqref{eq:istar} for every $S$
does \emph{not} in general force $\mathcal D\cap\mathbb R^L_{\ge0}=\{0\}$:
lying outside a convex cone is not preserved under nonnegative combinations, so
excluding every indicator $\one_S$ from $\mathcal D$ excludes no other
nonnegative vector by itself.  Concretely, a normal that is $+1$ on part of $P$
and $-1$ on part of $Q$ can satisfy $\nu\cdot\one_S<0$ while $\nu\cdot d\ge0$
for a nonnegative $d$ weighted towards $P$.  Everything below uses only the
implication as stated; where the cone condition itself is needed --- in the
constructions of Section~\ref{sec:suf} --- it is imposed directly.
\end{remark}

The case $S=\{i\}$ of \eqref{eq:istar} recovers Lemma~\ref{lem:prot} in the
form we need: for $i\in P$ there is a tight $T$ with $f\in T$ and $i\notin T$.

\begin{lemma}[exactness directions]\label{lem:edir}
There is $d^\ast\in\mathcal D$ with $\sum_{i\in P}d^\ast_i=\phi_f>0$, and
symmetrically one with $\sum_{i\in Q}d_i=\phi_g>0$.
\end{lemma}

\begin{proof}
Exactness at $\{f\}$ gives $z\in\core(v)$ with $z_f=v(\{f\})=0$;
Lemma~\ref{lem:rigid} gives $z(P)+z_f=v(S_1)=\phi(P)+\phi_f$, so
$\sum_{i\in P}(z_i-\phi_i)=\phi_f$.  Put $d^\ast=z-\phi$, which lies in
$\mathcal D$ because $z$ is a core point.
\end{proof}

When $p=1$ this pins the whole of $\phi_f$ onto a single poorer player.  When
$p\ge2$ it says only that some member of $P$ gains, and that slack is exactly
what the constructions of Section~\ref{sec:suf} exploit.

\begin{lemma}[crossing exclusion]\label{lem:xstar}
Let $T$ be $\phi$-tight.
\begin{enumerate}[nosep,label=(\roman*)]
\item If $f\in T$ then
 $s_\phi\bigl((P\cap T)\cup\{f\}\bigr)\le\phi(Q\setminus T)+[g\notin T]\,\phi_g$.
\item If $g\in T$ then
 $s_\phi\bigl((Q\cap T)\cup\{g\}\bigr)\le\phi(P\setminus T)+[f\notin T]\,\phi_f$.
\end{enumerate}
In particular, if $T\supseteq H$ and $P\cap T=\emptyset$ then
$\phi_f\le\phi(Q\setminus T)$, and since every $\phi_i$ with $i\in L$ is
strictly below $\phi_f$ this forces \eqref{eq:key}: $|Q\setminus T|\ge2$.
Symmetrically, if $T\supseteq H$ and $Q\cap T=\emptyset$ then
$|P\setminus T|\ge2$.
\end{lemma}

\begin{proof}
Cross $T$ with $S_1$ in Lemma~\ref{lem:cross}: here
$S_1\cap T=(P\cap T)\cup\{f\}$ and
$N\setminus(S_1\cup T)=(Q\setminus T)\cup(\{g\}\setminus T)$, which is (i);
(ii) is the mirror.  For the last part, $P\cap T=\emptyset$ makes the left side
$s_\phi(\{f\})=\phi_f$, and $g\in T$ deletes the $\phi_g$ term.  If
$|Q\setminus T|=0$ we get $\phi_f\le0$, false; if $|Q\setminus T|=1$ we get
$\phi_f\le\phi_i<\phi_f$, false.
\end{proof}

\section{Necessity}\label{sec:nec}

By the $f/g$ symmetry, which exchanges $(f,P,S_1)$ with $(g,Q,S_2)$ and carries
type $(p,q)$ to type $(q,p)$, it suffices to bound one side.

\begin{theorem}\label{thm:g1}
If $p=1$ then $q\ge3$.
\end{theorem}

\begin{proof}
Write $P=\{a\}$ and suppose $q\le2$.  By Lemma~\ref{lem:nonull} we may assume
$\phi_i>0$ on $L$, so no singleton is tight.

\emph{Step 1.}  Lemma~\ref{lem:improve} at $S=\{a\}$ requires a $\phi$-tight
$T$ with $|T\cap\{a\}|<\varepsilon_f$, hence $f\in T$ and $a\notin T$.

\emph{Step 2: $g\notin T$.}  If $g\in T$ then $T\supseteq H$ and
$P\cap T=\emptyset$, so \eqref{eq:key} gives $|Q\setminus T|\ge2$, forcing
$q=2$ and $T\cap Q=\emptyset$, i.e.\ $T\subseteq H$.  But $\{f\}$ is not tight
and $H$ is slack.  So $T\subseteq Q\cup\{f\}$, and $R:=T\cap Q\ne\emptyset$
because $\{f\}$ is not tight.

\emph{Step 3.}  Let $d^\ast$ be as in Lemma~\ref{lem:edir}; since $p=1$,
$d^\ast_a=\phi_f>0$.  The normal of $T=R\cup\{f\}$ has $\nu_a=-1$ and
$\nu_j=[j\in R]$ for $j\in Q$, so $\nu_T\cdot d^\ast\ge0$ reads
$\sum_{j\in R}d^\ast_j\ge d^\ast_a=\phi_f>0$; some $j\in R$ has
$d^\ast_j>0$.

\emph{Case $q=1$.}  Then $R=Q=\{b\}$ and $d^\ast_b>0$, so $d^\ast$ is strictly
positive on all of $L$: a nonzero nonnegative element of $\mathcal D$,
contradicting Lemma~\ref{lem:improve}.

\emph{Case $q=2$.}  Write $Q=\{b,c\}$ and, relabelling $b$ and $c$ if
necessary, take $d^\ast_b>0$.  If $d^\ast_c\ge0$ then again $d^\ast$
contradicts Lemma~\ref{lem:improve}, so $d^\ast_c<0$.  Apply
Lemma~\ref{lem:improve} at $S=\{a,b\}$: some $\phi$-tight $T'$ has
$|T'\cap\{a,b\}|<\varepsilon_f+\varepsilon_g$.  Enumerate by
$(\varepsilon_f,\varepsilon_g)$.
\begin{itemize}[nosep]
\item $(0,0)$: the right side is $0$, impossible.
\item $(1,0)$: $a,b\notin T'$ and $f\in T'$, so $T'\subseteq\{c,f\}$ and
      $T'\ne\{f\}$, whence $T'=\{c,f\}$, with $\nu_a=-1$, $\nu_b=0$,
      $\nu_c=1$.  Then $\nu\cdot d^\ast\ge0$ gives
      $d^\ast_c\ge d^\ast_a=\phi_f>0$, against $d^\ast_c<0$.
\item $(0,1)$: $a,b\notin T'$ and $g\in T'$, so $T'=\{c,g\}$, with $\nu_a=0$,
      $\nu_b=-1$, $\nu_c=0$.  Then $\nu\cdot d^\ast\ge0$ gives
      $d^\ast_b\le0$, against $d^\ast_b>0$.
\item $(1,1)$: $T'\supseteq H$ and $T'$ omits $a$ or $b$.  If $a\notin T'$,
      Step 2 already excluded it.  If $a\in T'$ and $b\notin T'$ then
      $\nu_a=0$, $\nu_b=-1$ and $\nu_c=[c\in T']-1$.  When $c\in T'$,
      $\nu\cdot d^\ast\ge0$ gives $d^\ast_b\le0$, again a contradiction; when
      $c\notin T'$ we have $T'=H\cup\{a\}$, and crossing that with $S_2$ gives
      $T'\cap S_2=\{g\}$ and $T'\cup S_2=N$, hence $\phi_g\le0$, false.
\end{itemize}
Every case is contradictory.
\end{proof}

\begin{theorem}[necessity]\label{thm:necessity}
A boundary partition of type $(p,q)$ requires $p,q\ge1$ and $p+q\ge4$.
\end{theorem}

\begin{proof}
$p,q\ge1$ was noted after Proposition~\ref{prop:descent}.
Theorem~\ref{thm:g1} gives $p=1\Rightarrow q\ge3$ and, by the $f/g$ symmetry,
$q=1\Rightarrow p\ge3$.  If $p,q\ge2$ then $p+q\ge4$ directly.  In all three
cases $p+q\ge4$.
\end{proof}

\section{Sufficiency: two closed-form families}\label{sec:suf}

Both families are lower envelopes, so exactness and normalization are free by
Lemma~\ref{lem:env}.  The leximin verification factors through a single lemma.

\begin{lemma}[pinning]\label{lem:pin}
Let $v$ be exact and singleton-normalized on $N$, let $c>1$, and put
$\phi=\one_L+c\,\one_H$, so $v(N)=m+2c$.  Suppose
\begin{enumerate}[nosep,label=(\alph*)]
\item $v(S_1)=\phi(S_1)$ and $v(S_2)=\phi(S_2)$, and these sum to $v(N)$;
\item the only $x\in\core(v)$ with $x\ge\one_N$ is $\phi$ itself.
\end{enumerate}
Then $\phi=\Lmin(v)$, its strata are exactly $L$ at level $1$ and then $H$ at
level $c$, the boundary is strict, and $S_1,S_2$ are $\phi$-tight.
\end{lemma}

\begin{proof}
By (a) and Lemma~\ref{lem:rigid}, $x(P)+x_f=p+c$ and $x(Q)+x_g=q+c$ throughout
the core.  Stage $1$ maximizes $\lambda$ such that some core point has all
coordinates at least $\lambda$; $\phi$ gives $\lambda\ge1$, and $\lambda>1$
would produce a core point strictly above $\one_N$, excluded by (b).  So
$\mu_1=1$, and the feasible set at that level is the single point $\phi$, whose
coordinates equal $1$ exactly on $L$; hence $B_1=L$.  At stage $2$ the players
of $L$ are fixed at $1$, so $x_f+x_g=(m+2c)-m=2c$, the largest common lower
bound is $c$, and it is attained only at $x_f=x_g=c$.  Strictness is $1<c$ and
tightness of $S_1,S_2$ is (a).
\end{proof}

Lemma~\ref{lem:pin} is Lemma~\ref{lem:improve} in constructive form: for a
two-stratum shape $L\prec H$ the first leximin stage is the only one that is
not automatic, so a construction need only control it.

\subsection{Family A: both sides large}

Let $p,q\ge2$ and $c=2$; the common total is $m+4$.  Order the coordinates
$x_1,\dots,x_p\mid y_1,\dots,y_q\mid f,g$, and write
\[
 \Pi=(2,2,1^{p-2}),\qquad \Sigma_a=(0,2,1^{q-2}),\qquad
 \Sigma_b=(0,4,1^{q-2}),
\]
which sum to $p+2$, $q$ and $q+2$ respectively, every entry of $\Pi$ being at
most $2$.  The family consists of
\[
\begin{array}{llll}
 A_0     &(\one^p            &\mid \one^q &\mid 2,2)\ =\ \phi,\\
 A_1     &(\Pi               &\mid \Sigma_a &\mid 0,2),\\
 A_2     &(\Pi               &\mid \Sigma_b &\mid 0,0),\\
 A_3^x   &(0\text{ at }x,\ 2\text{ at one other},\ 1\text{ elsewhere}
          &\mid \Sigma_a &\mid 2,2)\quad (x\in P),\\
 A_4^j   &(\Pi &\mid 0\text{ at }y_1,y_j,\ 1\text{ elsewhere} &\mid 0,4)
          \quad (j=2,\dots,q).
\end{array}
\]

\begin{proposition}\label{prop:famA}
For all $p,q\ge2$ the family above realizes a boundary partition of type
$(p,q)$, with $\Lmin(v)=\one_L+2\,\one_H$ and $s_\phi(H)=4$.
\end{proposition}

\begin{proof}
All entries are nonnegative and every row totals $m+4$.  Each coordinate
vanishes in some row --- $x\in P$ in $A_3^x$, $y_1$ in $A_1$, $y_j$ in $A_4^j$,
$f$ in $A_1$, $g$ in $A_2$ --- so Lemma~\ref{lem:env} gives an exact
singleton-normalized game.

\emph{Rigidity.}  Every row has $g(P)+g_f=p+2$ and $g(Q)+g_g=q+2$: for $A_0$
these are $p+2$ and $q+2$; for $A_1$, $(p+2)+0$ and $q+2$; for $A_2$,
$(p+2)+0$ and $(q+2)+0$; for $A_3^x$, $g(P)=0+2+(p-2)=p$ and $q+2$; for
$A_4^j$, $(p+2)+0$ and $(q-2)+4$.  Hence $v(S_1)=p+2$, $v(S_2)=q+2$, summing to
$v(N)$, which is hypothesis (a) of Lemma~\ref{lem:pin}.

\emph{Two pins.}  First, $\max_g(g_{y_1}+g_x)=2$ for each $x\in P$: row $A_0$
attains $2$, and every other row has $g_{y_1}=0$ with $g_x\le2$.  Hence
$v(N\setminus\{y_1,x\})=m+2$ and every core point obeys $x_{y_1}+x_x\le2$.
Second, every row has $g_g\ge g_f$ (namely $2\ge2$, $2\ge0$, $0\ge0$, $2\ge2$,
$4\ge0$) with equality at $A_0$, so $v(\{g\}\cup P)=p+2$ and every core point
obeys $x_g\ge x_f$.

\emph{Uniqueness above $\one_N$.}  Let $x\in\core(v)$ with $x\ge\one_N$.  The
first pin and $x_{y_1},x_x\ge1$ give $x_{y_1}=1$ and $x_x=1$ for all $x\in P$;
then rigidity gives $x_f=2$, the second pin gives $x_g\ge2$, and rigidity gives
$x(Q)=q+2-x_g\le q$.  But $x(Q)\ge q$, so $x(Q)=q$, every $x_{y_j}=1$ and
$x_g=2$.  Hence $x=\phi$, which is hypothesis (b).

\emph{Slack.}  $v(H)=\min_g(g_f+g_g)=0$, attained at $A_2$, while $\phi(H)=4$.

Lemma~\ref{lem:pin} now applies with $c=2$, and $S_1,S_2$ are tight, disjoint,
covering, with traces $\{f\}$ and $\{g\}$.
\end{proof}

The two hypotheses $p\ge2$ and $q\ge2$ are visible in the construction.  $\Pi$
must fit $p+2$ into $p$ entries each at most $2$, which needs $p+2\le2p$; and
that ceiling of $2$ is exactly the first pin.  At $q=1$ the row normalizing $g$
would need $g_{y_1}=q+2=3$ against the same ceiling, forcing $g(P)=0$ against
$g(P)=p+2$.  These are the two readings of \eqref{eq:key}.

\subsection{Family B: one side a singleton}

Let $p=1$, $q\ge3$, and $c=q-1\ge2$; the common total is $3q-1$.  Order the
coordinates $a\mid y_1,\dots,y_q\mid f,g$.  Four rows suffice, for every $q$:
\[
\begin{array}{llll}
 B_0 &(1     &\mid \one^q                &\mid c,\ c)\ =\ \phi,\\
 B_1 &(0     &\mid q,\ q-1,\ 0^{q-2}     &\mid 1+c,\ 0),\\
 B_2 &(1+c   &\mid 0,\ q-1,\ 0^{q-2}     &\mid 0,\ q),\\
 B_3 &(1     &\mid 1,\ 0,\ 1^{q-2}       &\mid c,\ q).
\end{array}
\]

\begin{proposition}\label{prop:famB}
For all $q\ge3$ the family above realizes a boundary partition of type
$(1,q)$, with $\Lmin(v)=\one_L+(q-1)\one_H$ and $s_\phi(H)=q-2$.
\end{proposition}

\begin{proof}
Every row totals $3q-1$, and coordinates vanish: $a$ in $B_1$, $y_1$ in $B_2$,
$y_2$ in $B_3$, $y_k$ ($k\ge3$) in $B_1$, $f$ in $B_2$, $g$ in $B_1$.  So
Lemma~\ref{lem:env} applies.

\emph{Rigidity.}  Every row has $g_a+g_f=1+c=q$ and $g(Q)+g_g=q+c=2q-1$; this
is checked row by row.  Hence hypothesis (a) of Lemma~\ref{lem:pin}.

\emph{Two pins.}  Using rigidity, tightness of $\{a,y_1\}\cup H$ is equivalent
to $\max_g\,g(Q\setminus\{y_1\})=q-1$, which holds: the four rows give $q-1$,
$q-1$, $q-1$ and $q-2$.  So every core point obeys $x(Q\setminus\{y_1\})\le
q-1$.  Likewise tightness of $\{y_2\}\cup H$ is equivalent to
$\max_g\,(g_a+g(Q\setminus\{y_2\}))=q$, which holds: the four rows give
$1+(q-1)$, $0+q$, $(1+c)+0$ and $1+1+(q-2)$.  So every core point obeys
$x_a+x(Q\setminus\{y_2\})\le q$.

\emph{Uniqueness above $\one_N$.}  Let $x\in\core(v)$ with $x\ge\one_N$.  The
first pin bounds a sum of $q-1$ terms each at least $1$ by $q-1$, so
$x_{y_j}=1$ for $j\ge2$.  The second bounds a sum of $q$ terms each at least
$1$ by $q$, so $x_a=1$ and $x_{y_j}=1$ for $j\ne2$; in particular
$x_{y_1}=1$.  Rigidity then gives $x_f=(1+c)-1=c$ and $x_g=(q+c)-q=c$, so
$x=\phi$.

\emph{Slack.}  $v(H)=\min_g(g_f+g_g)=\min\{2c,\ q,\ q,\ c+q\}=q$, while
$\phi(H)=2c=2q-2$; the slack is $q-2$, positive exactly for $q\ge3$.

Lemma~\ref{lem:pin} applies with $c=q-1$.
\end{proof}

At $q=2$ the same four rows still define an exact game, but two things fail at
once: the slack $q-2$ vanishes, and $c=q-1=1$ makes $\phi$ constant, so $H$ is
not a strict richer suffix.  Instantiating them gives rows $(1,1,1,1,1)$,
$(0,2,1,2,0)$, $(2,0,1,0,2)$, $(1,1,0,1,2)$ with $\Lmin=(1,1,1,1,1)$ and
$v(H)=\phi(H)=2$.  This is the constructive shadow of Theorem~\ref{thm:g1}: at
$q=2$ the coalition $\{y_2\}\cup H$ used in the second pin would have to
satisfy \eqref{eq:key} with $|Q\setminus T|=1$.

\begin{theorem}[sufficiency]\label{thm:sufficiency}
For every $p,q\ge1$ with $p+q\ge4$ there is an explicit exact
singleton-normalized game carrying a boundary partition of type $(p,q)$.
\end{theorem}

\begin{proof}
If $p,q\ge2$ use Family A.  If $p=1$ then $p+q\ge4$ forces $q\ge3$, so Family B
applies; if $q=1$ then $p\ge3$ and the mirror of Family B applies.  The three
ranges are disjoint and exhaust the hypothesis.
\end{proof}

Theorems~\ref{thm:necessity} and~\ref{thm:sufficiency} together are
Theorem~\ref{thm:main}.

\section{At the threshold}\label{sec:frontier}

Step 2 of Theorem~\ref{thm:g1} shows that when $q\le2$ no $\phi$-tight
coalition contains $H$ and omits $P$.  At $q=3$ the opposite holds.

\begin{theorem}[frontier dichotomy]\label{thm:g4}
In every boundary partition of type $(1,3)$ there is a $\phi$-tight $T$ with
$H\subseteq T$ and $a\notin T$; necessarily $T=H\cup\{j\}$ for a single
$j\in Q$.
\end{theorem}

The shape of $T$ is forced by \eqref{eq:key}, which requires
$|Q\setminus T|\ge2$, together with slackness of $H$, which excludes
$|Q\setminus T|=3$.  Existence is established by an exhaustive finite check
over the payoff stratifications, described in Section~\ref{sec:status}; we have
no human proof of it.  Together with Step 2 of Theorem~\ref{thm:g1} the
dichotomy reads:
\[
 q\le2:\ \text{such a }T\text{ is impossible};\qquad
 q=3:\ \text{such a }T\text{ is mandatory}.
\]
That is the sense in which the threshold is sharp.  The object the obstruction
forbids does not become gradually easier to build as $q$ grows; it passes from
unavailable to unavoidable at a single step.

\begin{remark}[a claim we withdraw]\label{rem:withdraw}
The crossing budget behind \eqref{eq:key} reads
$\phi_f\le\phi(Q\setminus T)$, and in the first type-$(1,3)$ witness we
harvested it held with equality.  We had recorded that as evidence that the
threshold is sharp in a second, quantitative sense.  It is not: the budget is
in general strict.  A search over all stratifications returns $74$
stratification--$j$ pairs, each carrying an audited type-$(1,3)$ realization
in which it is strict ---
for instance
$\phi=(25/24,1,49/48,17/16,13/12,13/12)$ on $(a,b,c,d,f,g)$, with
$\max_L\phi=17/16<13/12=\min_H\phi$, $T=\{c\}\cup H$ tight, and
$\phi_f=13/12<33/16=\phi_b+\phi_d$, a gap of $47/48$.  Moreover the same search
with the strictness dropped is satisfiable on exactly the same $74$
stratification--$j$ pairs, so strictness costs nothing.  Theorem~\ref{thm:g4}
is the correct sharpness statement and the equality claim is withdrawn.
\end{remark}

\section{The seven-player consequence}\label{sec:zp}

We can now settle the zero-payoff question that motivated
Definition~\ref{def:bp}.  A census of seven-player games had left one
admissible profile undecided, namely $\mathrm{ZP}(2,2)$ with five poorer
players; by Proposition~\ref{prop:descent}(v) it descends either to a boundary
partition of type $(2,2)$ or to $\mathrm{ZP}(1,2)$ or $\mathrm{ZP}(2,1)$, and
those descend in turn to boundary partitions of type $(1,2)$ or $(2,1)$.  We
had conjectured all three types impossible, which would have refuted the
profile.  Theorem~\ref{thm:main} says otherwise: types $(1,2)$ and $(2,1)$ are
indeed impossible, but $(2,2)$ is not.

\begin{example}[a type-$(2,2)$ game]\label{ex:22}
On $N=\{a,\dots,f\}$ take the lower envelope of the five allocations of common
total $8$
\[
 (0,0,1,4,2,1),\quad(0,0,4,4,0,0),\quad(4,0,0,0,2,2),\quad
 (1,0,2,3,0,2),\quad(0,1,2,3,1,1).
\]
Then $v$ is exact and singleton-normalized, $\Lmin(v)=(1,1,2,2,1,1)$ with
strata $\{a,b,e,f\}\prec\{c,d\}$, so $H=\{c,d\}$ is a strict richer suffix with
$m=4$; $v(cd)=0$, so the boundary is slack with $s_\phi(H)=4$; the proper
tight coalitions are $ad$, $abd$, $ade$, $acde$, $adf$, $acdf$, $cef$, $def$,
$acdef$; and $S_1=cef$, $S_2=abd$ are tight, disjoint, cover $N$, and meet $H$
in $\{c\}$ and $\{d\}$.  Their poorer parts are $\{e,f\}$ and $\{a,b\}$, so the
type is $(2,2)$.
\end{example}

\begin{corollary}\label{cor:zp7}
There is a seven-player exact singleton-normalized game carrying a
$\mathrm{ZP}(2,2)$ configuration with five poorer players.
\end{corollary}

\begin{proof}
Adjoin a null player $g$ to Example~\ref{ex:22} --- concretely, append a
seventh coordinate $0$ to each of the five rows.  By
Proposition~\ref{prop:descent}(iv) the result is exact and singleton-normalized
with $\Lmin=(1,1,2,2,1,1,0)$, strata $\{g\}\prec\{a,b,e,f\}\prec\{c,d\}$, the
same boundary $H=\{c,d\}$ with slack $4$, and tight sets $cefg$ and $abdg$
meeting in $\{g\}$ and covering $N$; the poorer parts are $\{e,f\}$ and
$\{a,b\}$.
\end{proof}

Family A of Section~\ref{sec:suf} gives an alternative, uniform witness at
$(2,2)$, with $\Lmin=\one_L+2\one_H$; Example~\ref{ex:22} is retained because
it is the game that refuted the conjecture and because its tight family is
visibly different, which is a useful reminder that the families of
Section~\ref{sec:suf} are constructions and not normal forms.

\begin{remark}[why the census missed it]
Type $(2,2)$ was absent from a census of some $20{,}000$ sampled slack
boundaries with $m=4$; a later targeted harvest found it once in $15{,}967$
games, against six of type $(1,3)$ and seven of type $(3,1)$.  It is rare, not
absent.  The general lesson --- that ``configuration $X$ never occurs'' is a
statement about a sampler until it is a statement about set systems --- is
recorded for the same reason in \citep{Flanagan2026phase}.
\end{remark}

\section{What is proved, what is checked, what is verified}\label{sec:status}

\begin{center}
\begin{tabular}{lp{0.62\textwidth}}
\toprule
\textbf{status} & \textbf{statement}\\
\midrule
proved &
 Theorems~\ref{thm:exact} and~\ref{thm:farkas};
 Lemmas~\ref{lem:prot}, \ref{lem:rigid}, \ref{lem:normal}, \ref{lem:improve},
 \ref{lem:edir}, \ref{lem:xstar}, \ref{lem:nonull}, \ref{lem:pin};
 Proposition~\ref{prop:descent}; Theorems~\ref{thm:g1}
 and~\ref{thm:necessity}; Propositions~\ref{prop:famA}
 and~\ref{prop:famB} and Theorem~\ref{thm:sufficiency}; hence
 Theorem~\ref{thm:main}.  None of these consults a machine.\\
finite check &
 Theorem~\ref{thm:g4}: unsatisfiable over all $60$ payoff stratifications of
 type $(1,3)$, using the encoding of Theorems~\ref{thm:exact}
 and~\ref{thm:farkas}, with a satisfiable control confirming the encoding is
 not vacuous.\\
verified &
 Example~\ref{ex:22} and Corollary~\ref{cor:zp7}, by exact rational
 recomputation of the envelope, of exactness at definition level, of $\Lmin$,
 of the tight family and of the configuration; the $74$ strict-budget
 realizations (one per stratification--$j$ pair) of
 Remark~\ref{rem:withdraw}, each audited the same way; and
 the two families of Section~\ref{sec:suf} at all $18$ instances with
 $n\le9$.\\
\bottomrule
\end{tabular}
\end{center}

An earlier route to Theorem~\ref{thm:necessity} was a finite check --- $48$
unsatisfiable instances at $m=3$ over the encoding of
Section~\ref{sec:methods}, cross-checked by an independent enumeration of tight
families --- and it is retained in the repository as a control on the hand
proof rather than as a dependency.  Likewise the padding construction that
first established sufficiency (adjoin a null player and place it on either side
of the partition) remains true but is no longer load-bearing, having been
superseded by Section~\ref{sec:suf}.

\section{Scope}\label{sec:scope}

Four limitations are worth stating plainly.

First, the characterization is specific to $|H|=2$.  For $|H|\ge3$ a
``boundary partition'' would have to be defined with more than two blocks, and
Lemma~\ref{lem:rigid} would still apply, but neither the necessity argument nor
the families have been examined in that setting.

Second, Theorem~\ref{thm:g4} is a finite check, not a theorem with a human
proof.  It is not used anywhere else in this note; nothing in
Theorem~\ref{thm:main} depends on it.

Third, the families of Section~\ref{sec:suf} carry no minimality claim.
Family A uses $m+2$ rows and Family B uses four, but we prove no lower bound on
the number of generators, and Family A's uniform shape --- payoffs $1$ and $2$,
slack $4$ independent of $p$ and $q$ --- is a property of that construction
rather than of the configuration.  Example~\ref{ex:22} and the configurations of
Remark~\ref{rem:withdraw} look nothing like it.

Fourth, we characterize which types occur, not how often or in what
proportion of the exact cone.  The census figures quoted in
Section~\ref{sec:zp} are descriptive of a sampler.

\section*{Reproducibility note}

Theorem~\ref{thm:main} and everything it rests on are proved by hand; the
families of Section~\ref{sec:suf} are finite verifications of row arithmetic,
uniform in $(p,q)$.  The one proof dependency that is a machine result is
Theorem~\ref{thm:g4}, and it is used nowhere else.  Exact rational programs
archived in \texttt{../verifiers/part} reconstruct, independently of the text:
the game of Example~\ref{ex:22} and its seven-player suspension, including
exactness at definition level by linear programming over the core rather than
by vertex enumeration; the two families at every instance with $n\le9$; the
finite check behind Theorem~\ref{thm:g4} together with its control; the
withdrawn-equality search of Remark~\ref{rem:withdraw}; and the superseded
$m=3$ finite check with its independent cross-check.  All arithmetic is exact;
every solver call is asserted satisfiable or unsatisfiable, and no call
returned an indeterminate result.  No core-vertex enumeration is performed at
any $n$.

% NOTE (lead, 2026-08-08): bibliography pruned to entries that are (a) cited
% in the text and (b) verified against the companion papers' bibliographies.
% A prior-art positioning pass (with web-verified metadata, as done for the
% companion papers) is PENDING before this note goes anywhere; expected
% additions include the egalitarian-selection literature cited in the
% companions (Arin--I\~narra, Dutta--Ray, Llerena--Mauri, Dietzenbacher) once
% the positioning paragraphs that cite them are written.

\noindent\emph{Draft note on references.}  This list contains only the works
actually cited above.  Every entry is carried over verbatim from the companion
paper's bibliography and still requires the editorial verification pass; no
reference here was reconstructed from memory.  See the prior-art paragraph in
Section~\ref{sec:intro-priorart} for the citations this note still needs.

\begin{thebibliography}{99}

\bibitem[Flanagan(2026)]{Flanagan2026phase}
Flanagan M (2026) Balanced transport at the egalitarian core: a complete
boundary phase diagram. Companion manuscript.

\bibitem[Lohmann et al.(2012)]{LohmannBormHerings2012}
Lohmann E, Borm P, Herings PJJ (2012) Minimal exact balancedness.
Mathematical Social Sciences 64:127--135.

\bibitem[Studen\'y and Kratochv\'il(2022)]{StudenyKratochvil2022}
Studen\'y M, Kratochv\'il V (2022) Facets of the cone of exact games.
Mathematical Methods of Operations Research 95:35--80.

\end{thebibliography}

\end{document}
